\subsection{\texorpdfstring{\(\mathbf{B}_{\mathrm{M}}\) 为正且退化的情形}{B\_M 为正且退化的情形}}

在本号中，我们假设 S 有限，\((W, S)\) 不可约，且型 \(B_{M}\) 是正的且退化的。

引理 2. 正交补 \(\mathbf{E}^0\) —— 即 \(\mathbf{E}\) 关于 \(\mathbf{B}_{\mathbf{M}}\) 的正交补 —— 是 1 维的；它由一个元素 \(v = \sum_{s\in S}v_s e_s\) （其中 \(v_{s} > 0\) 对所有 \(s\) 成立）生成。

这由 §3 第 5 号的引理 4（应用于以 \(\mathrm{B}_{\mathbb{M}}(e_s, e_{s'})\) 为元素的矩阵）得出。

设 \(v = \sum_{s} v_s e_s\) 是满足引理 2 的条件且使 \(\sum_{s} v_s = 1\) 的向量，并设 A 是 \(E^*\) 中由 \(y^* \in E^*\) （满足 \(\langle v, y^* \rangle = 1\) ）构成的仿射超平面。若 T 表示 \(v\) 在 \(E^*\) 中的正交补，则 A 具有一个仿射空间的自然结构，以 T 为平移空间。此外，型 \(B_M\) 通过过渡到商在 \(E/E^0\) 上定义一个非退化标量积，从而也在其对偶 T 上定义；这给出仿射空间 A 上的一个欧氏结构（《代数》第 IX 章，§6，第 6 号）。

设 G 是 \(\mathbf{GL}(\mathbf{E})\) 的由保持 v 与 \(B_{M}\) 不变的自同构构成的子群；若 \(g \in G\) ，则逆步映射 \({}^{t}g^{-1}\) 保持 A 与 T 稳定，且通过限制到 A 上定义 A 的一个位移 \(i(g)\) （参见~§3）。立即可见这给出从 G 到 A 的位移群的一个同构。此外，A 的点 a 的稳定子 \(G_{a}\) 可与 Hilbert 空间 T 的正交群等同，从而是紧的。另一方面，G 是一个在无穷远处可数的局部紧群，而 A 是一个 Baire 空间：由此（《积分》第 VII 章，附录，引理 2）可知映射 \(\psi : g \mapsto g(a)\) 定义从 \(G/G_a\) 到 A 的一个同胚。于是 G 适当地作用于 A（《一般拓扑》第 III 章，§4，第 2 号，命题 5）。由于 W 是 G 的子群，它可与 A 的一个位移群等同。我们将证明这个群满足 §3 的假设。更确切地说：

命题 10. 赋予离散拓扑的群 W 适当地作用于 A；它由正交反射生成；它是无限的、不可约的且本质的（§3，第 7 号）。交 C ∩ A 是 W 在 A 中的一个室。若 L\_s 表示 A 中由 A 与 E* 的正交于 e\_s 的超平面相交而得的超平面，则诸 L\_s （对于 s ∈ S）是 C ∩ A 的墙。若 ε\_s 是 T 中正交于 L\_s 且与 C ∩ A 位于 L\_s 同侧的单位向量，则 (ε\_s\textbar ε\_t) = -cos(π/m(s,t)) （对于 s,t ∈ S），且 W 的 Coxeter 矩阵（§3，第 4 号）与 M 等同。

由定理 1 的推论 3，W 在 \(\mathbf{GL}(\mathbf{E})\) 中离散，从而在 \(\mathbf{G}\) 中离散，且适当地作用于 A。设 \(s \in S\) 。由于 \(\operatorname{Card} S \geqslant 2\) ，\(\mathbf{E}^*\) 中正交于 \(e_s\) 的超平面不正交于 \(v\) ，它与 A 的交 \(\mathbf{L}_s\) 确是超平面。于是与 \(s\) 相应的位移是保持 \(\mathbf{L}_s\) 的所有点固定的 2 阶位移：它必然是与 \(\mathbf{L}_s\) 相伴随的正交反射。由此可知 W 由正交反射生成。定理 2 现在表明它是无限的，而命题 7 表明它是本质的且不可约的。

由于 C 是一个开单形锥，其墙是方程为 \(\langle x^{*}, e_{s} \rangle = 0\) （对于 \(s \in S\) ）的超平面，交 \(C \cap A\) 是 A 中诸 \(L_{s}\) 之并的补集的一个凸的、从而连通的开且闭子集。此外，\(C \cap A\) 非空，因为若 \(x^{*} \in C\) ，我们有 \(\langle x^{*}, v \rangle = \sum_{s} v_{s} \langle x^{*}, e_{s} \rangle > 0\) 且 \(\langle x^{*}, v \rangle^{-1} x^{*} \in C \cap A\) 。由此可知 \(C \cap A\) 是 A 相对于诸 \(L_{s}\) 的系统的一个室。此外，\(w(C \cap A) \cap L_{s} = \varnothing\) （对所有 \(w \in W\) ；参见~第 4 号，性质 \((\mathrm{P}_{n})\) ）成立，因而 \(C \cap A\) 是 A 相对于由诸 \(L_{s}\) 经 W 的元素变换所得的系统的一个室；由 §3 第 2 号定理 1 的推论可知，\(C \cap A\) 是 A 相对于 W 的一个室。

设 \(a_{s}^{*}\) 是单形 \(C \cap A\) 的不在 \(L_{s}\) 中的顶点。我们有

\[
\left\langle a _ {s} ^ {*}, e _ {t} \right\rangle = 0
\]

对 \(s,t\in \mathrm{S}\) 且 \(s\neq t\) 成立，且

\[
\langle a _ {s} ^ {*}, e _ {s} \rangle = v _ {s} ^ {- 1} \langle a _ {s} ^ {*}, v \rangle = v _ {s} ^ {- 1}.
\]

设 \(\varepsilon_{s}\) 是 T 中由下列关系式定义的向量：

\[
\left(\varepsilon_ {s} \mid a _ {s} ^ {*} - a _ {t} ^ {*}\right) = v _ {s} ^ {- 1} \text {for} t \in S, t \neq s.
\]

向量 \(\varepsilon_{s}\) 正交于 \(L_{s}\) ，且关于超平面 \(L_{s}\) 与 \(C \cap A\) 位于同侧。此外

\[
\left(\varepsilon_ {s} \mid a _ {s} ^ {*} - a _ {t} ^ {*}\right) = \left\langle e _ {s}, a _ {s} ^ {*} - a _ {t} ^ {*} \right\rangle \text {for all} s, t \in S
\]

这表明 \(\varepsilon_{s}\) 是 \(e_{s}\) 的类在从 \(E/E^{0}\) 到 T 的同构（由二次型 \(B_{M}\) 给出）下的像。由此得到

\[
(\varepsilon_ {s} | \varepsilon_ {t}) = \mathrm{B} _ {\mathrm{M}} (e _ {s}, e _ {t}).
\]

因此 \(\varepsilon_{s}\) 是单位向量，而命题 10 的最后一个断言得证。

配备群 W 的欧氏仿射空间 A 称为与 Coxeter 矩阵 M 相伴随的空间，我们记之为 \(A_{M}\) 。

命题 10 有一个逆命题：

命题 11. 设 W 是欧氏仿射空间 A 的一个位移群，满足 §3 的假设。假设 W 是无限的、本质的且不可约的。则与 W 的 Coxeter 矩阵 M 相伴随的型 \(B_{M}\) 是正的且退化的，并且存在唯一的从与 M 相伴随的仿射空间 \(A_{M}\) 到 A 的同构，它与 W 的作用交换。这个同构把 \(A_{M}\) 的标量积变为 A 的标量积的倍数。

设 \(C_0\) 是 A 的一个室，并设 S 是相对于 \(C_0\) 的墙的正交反射所成的集合。若 \(\eta_s\) 表示正交于与反射 s 相伴随的超平面 \(N_s\) 且关于 \(N_s\) 与 \(C_0\) 位于同侧的单位向量（§3，命题 3），则型 \(B_M\) 满足 \(B_M(e_s, e_t) = (\eta_s | \eta_t)\) （对于 \(s, t \in S\) ）。从而它是正的。由于诸 \(\eta_s\) 线性相关（§3，第 9 号，命题 8），它是退化的。

于是我们可以把前面的构造应用于 M。沿用与上面相同的记号，\((\varepsilon_{s}|\varepsilon_{t}) = (\eta_{s}|\eta_{t})\) ，且存在唯一的 Hilbert 空间同构 \(\varphi\) ，从 T 到 A 的平移空间，使得 \(\varphi(\varepsilon_{s}) = \eta_{s}\) 。设 a 与 b 是 \(C_{0}\) 的两个不同顶点，\(s_{0}\) 是 S 中使 \(a \notin N_{s_{0}}\) 的反射。置 \(\lambda = (\eta_{s}|a - b)\) ，并设 \(\psi\) 是由下式定义的从 \(A_{M}\) 到 A 的仿射双射

\[
\psi \left(a _ {s _ {0}} + x\right) = a + v _ {s _ {0}} \lambda \varphi (x) \text {for} x \in \mathrm{T}.
\]

于是立即可见 \(\psi(\mathrm{L}_{s}) = \mathrm{N}_{s}\) （对所有 \(s \in S\) ），且 \(\psi\) 把 \(A_{M}\) 的标量积变为 A 上标量积的倍数。立即得到 \(\psi\) 与 W 的作用交换。最后，\(\psi\) 的唯一性是明显的，因为例如 \(a_{s}\) 是 \(A_{M}\) 中在诸反射 \(t \in S, t \neq s\) 下不变的唯一点。

